% This section is the paper's spine: it is what makes the work an adjudication rather than
% another entry in the debate. Each subsection states the explanation, who proposed it, and
% the operator that removes it.

\begin{table}[t]
\centering\small
\begin{tabular}{llp{6.2cm}}
\toprule
 & Explanation & Removed by \\
\midrule
C1 & Support shrinkage (sharpening) & --- (the residual) \\
C2 & Overtraining on saturated problems & Restrict to $p_{\text{base}} \in (0.05, 0.80)$ \\
C3 & Contamination of the base's coverage & Restrict to post-release problems \\
C4 & Answer-only scoring rewards guesses & Null-gold false-positive correction \\
C5 & Truncation (4k base vs.\ 32k descendants) & Retro-truncate both arms to a common budget \\
C6 & Matched-$k$ is not matched-compute & Compare at equal total sampled tokens \\
C7 & Prompt/template mismatch & Score both arms on the shared variant \\
\bottomrule
\end{tabular}
\caption{The seven live explanations. C3, C5, C6 and C7 have not previously been isolated.}
\label{tab:confounds}
\end{table}

\todo{one short paragraph per confound}
